\section{Introduction}

Metric Diophantine approximation along a fixed integer sequence $(q_n)$ asks
how the distances $\|q_nx\|$ to the nearest integer behave for typical and for
exceptional real numbers $x$. When $(q_n)$ is the sequence of convergent
denominators of an irrational $\alpha$, a criterion of Kraaikamp and
Liardet~\cite{KL} identifies the group $\Z+\Z\alpha$ through the size of
$\|q_nx\|$ relative to the approximation error of $\alpha$. Beyond such
zero-level statements, sets defined by a condition that is required for
\emph{all large} $n$ (rather than for infinitely many $n$) are much less
understood than their limsup counterparts; see for instance
Bugeaud, Liao and Rams~\cite{BLR} and the discussion in
Section~\ref{sec:prior}. In this paper we treat one sequence for which the
whole picture can be determined: a sparse, super-exponentially growing
subsequence of the convergent denominators of $\ee$. For this sequence we
compute the full weighted approximation spectrum, the Hausdorff dimension of
every level set, and the exact Hausdorff gauge of every level.

\subsection{Setup}

The continued fraction of $\ee$ is $[2;1,2,1,1,4,1,1,6,\dots]$. Let $P_j/Q_j$
be the convergent immediately preceding the partial quotient $2j$; the first
few are $3/1,19/7,193/71,2721/1001$. The denominators satisfy
\begin{equation}\label{eq:rec}
Q_{j+1}=(4j+2)Q_j+Q_{j-1}\quad(j\ge2),\qquad Q_1=1,\ Q_2=7,
\end{equation}
so $Q_{j+1}/Q_j\sim 4j$ and $\log Q_j\sim j\log j$ (super-exponential growth).
Write $\|y\|=\dist(y,\Z)$ and define
\[
\Lam(x):=\limsup_{j\to\infty} j\,\|Q_j x\|\in[0,\infty].
\]
More generally, for $\alpha\ge0$ put
$\Lam_\alpha(x):=\limsup_{j\to\infty} j^\alpha\|Q_j x\|\in[0,\infty]$,
so that $\Lam=\Lam_1$, and write $E_c:=\{x:\Lam(x)\le c\}$.
Let $L:=\Z+\Z\frac{\ee+1}{2}$, a countable dense subgroup of $\R$.

\subsection{Main results}

\paragraph{The spectrum.} The function $\Lam$ vanishes exactly on $L$, and
$\Lam\ge1/4$ off $L$ (Theorems~\ref{thm:class} and~\ref{thm:gap}). The
whole spectrum is quantized (Theorem~\ref{thm:full-spectrum}):
\[
\Lam(\R)=\frac14\Z_{\ge0}\cup\{\infty\}.
\]
Thus every interval $(k/4,(k+1)/4)$ is a gap of the spectrum, and the sublevel
sets are constant across each gap: $E_c=E_{k/4}$ for $c\in[k/4,(k+1)/4)$
(Corollary~\ref{thm:sub}).

\paragraph{Dimension.} Every finite sublevel set $E_c$ has Hausdorff dimension
$0$ (Theorem~\ref{thm:finite}), while $\{\Lam=\infty\}$ has full Lebesgue
measure (Corollary~\ref{cor:infinite-level}). For the subcritical weights
$0\le\alpha<1$, both $\{x:\Lam_\alpha(x)\le c\}$ and the prescribed-limit sets
$\{x\in[0,1]:j^\alpha\|Q_jx\|\to c\}$ have Hausdorff dimension exactly
$1-\alpha$ (Theorems~\ref{thm:upper} and~\ref{conj:weight}). The weight $j$ is
therefore critical: it is the threshold at which finite levels collapse to
dimension $0$.

\paragraph{Exact gauges.} Dimension does not distinguish the levels $k/4$, but
Hausdorff gauges do. For $k\ge1$ let $h_k$ be the gauge with
$h_k(1/Q_n)=(2k+1)^{-n}$, interpolated linearly between these scales. On the
circle $\R/\Z$, the level $\{\Lam=k/4\}$ has infinite but $\sigma$-finite
$\mathcal H^{h_k}$-measure, and the strict sublevel $\{\Lam<k/4\}$ is
$\mathcal H^{h_k}$-null (Theorem~\ref{thm:gauge}). The gauge $h_k$ is exact:
the $\mathcal H^{h_k}$-measure of the level is infinite in every nonempty open
arc, smaller gauges annihilate $\{\Lam\le k/4\}$, and larger gauges destroy
$\sigma$-finiteness (Corollary~\ref{cor:exact-gauge}). More generally, every
nondecreasing gauge is classified by the lower limit of its normalized values
along the scales $1/Q_n$, also when these values oscillate
(Corollary~\ref{cor:liminf-gauges}).

\paragraph{A trichotomy at each level.} Fix $k\ge1$ (Theorem~\ref{thm:trichotomy}).
\begin{enumerate}[label=(\roman*),leftmargin=2em]
\item The limsup level $\{\Lam=k/4\}$ has critical gauge base $2k+1$, as above.
\item The exact-limit set $\{x:j\|Q_jx\|\to k/4\}$ has critical gauge base $2$,
independently of $k$.
\item Each signed exact-limit set
$\{x:j\,(Q_jx-\mathrm{round}(Q_jx))\to\pm k/4\}$ is a single coset of $L$, hence
countable.
\end{enumerate}

\subsection{Method}

Everything rests on an exact symbolic coding of the level sets. For $x\in\R$
let $M_j$ be a nearest integer to $Q_jx$ and $u_j=Q_jx-M_j$. Because
$(Q_j)$ satisfies the integer recurrence~\eqref{eq:rec}, the numbers
\[
d_j:=M_{j+1}-(4j+2)M_j-M_{j-1}=(4j+2)u_j+u_{j-1}-u_{j+1}
\]
are integers, and $d_j-4ju_j\to0$ whenever $\Lam(x)<\infty$. Hence
$4\Lam(x)=\limsup_j|d_j|$, which gives the quantization of the spectrum, and the
level $\{\Lam\le k/4\}$ is described by digit sequences that eventually take
values in $\{-k,\dots,k\}$. Conversely, every such digit sequence is realized by
a point $x=\lim_jM_j/Q_j$ defined through a telescoping series
(Section~\ref{sec:gauge}). The level sets are thus Moran-type sets with
$2k+1$ branches per step at scale $1/Q_n$. The covering and
mass-distribution arguments that turn such a coding into gauge statements are
standard for homogeneous Moran sets (cf.\ Feng, Wen and Wu~\cite{FWW} and
Cabrelli, Mendivil, Molter and Shonkwiler~\cite{CMMS}). The specific inputs
here are the coding itself, the uniform branching bound that the critical
weight $j$ provides (which gives dimension $0$ for all finite levels,
Theorem~\ref{thm:finite}), and translation invariance under the dense group
$L$, which upgrades positive measure to infinite measure in every arc.

\subsection{Relation to prior work}\label{sec:prior}

The closest prior result is the criterion of Kraaikamp and Liardet \cite{KL},
quoted as Lemma~4.1 in \cite{BDS}. Let $\alpha$ be irrational with convergents
$p_n/q_n$. Then $\beta\notin\Z+\Z\alpha$ if and only if
$\|q_n\beta\|>\tfrac14\,q_n|q_n\alpha-p_n|$ for infinitely many $n$. This is a
classification together with a gap, with constant $1/4$, for the normalized
quantity $\|q_n\beta\|/(q_n|q_n\alpha-p_n|)$ along \emph{all} convergent
denominators. Theorems~\ref{thm:class} and~\ref{thm:gap} compare with it in two
ways.
\begin{itemize}[leftmargin=1.8em]
\item \emph{As a subsequence of the convergents of $\ee$.} Passing to every
third convergent denominator of $\ee$ enlarges the zero group from $\Z+\Z\ee$
to $L$: for $x=(\ee+1)/2$ the quantity $j\|Q_jx\|$ tends to $0$. It also
changes the normalization. Since $2j+1<1/(Q_j|Q_j\ee-P_j|)<2j+2$, one has
$\limsup_{j\to\infty}\|Q_jx\|/(Q_j|Q_j\ee-P_j|)=2\Lam(x)$, so the gap
$\Lam\ge1/4$ reads $\limsup\ge1/2$ in their normalization; the value
$\Lam(x_0)=1/4$ is attained by the point $x_0$ of Remark~\ref{rem:unweighted}.
\item \emph{As the full convergent sequence of $\beta=(\ee+1)/2$.} With
$R_j=(P_j+Q_j)/2\in\Z$, the fractions $R_j/Q_j$ are the convergents of
$\beta=[1;1,6,10,14,\dots]$. Applying the criterion of \cite{KL} (in the form
on p.~106 of \cite{BDS}) to $\beta$ gives $\|Q_jx\|>\tfrac14Q_j|Q_j\beta-R_j|$
infinitely often for $x\notin L$. Since $jQ_j|Q_j\beta-R_j|\to1/4$, this
yields $\Lam(x)\ge1/16$ off $L$. It already recovers $\{\Lam=0\}=L$; the direct
recurrence analysis gives the sharp gap $1/4$ and the complete quantization.
\end{itemize}
We therefore regard Theorems~\ref{thm:class} and~\ref{thm:gap} as a variant of
known results. The new results are the complete quantization of the spectrum
(Theorem~\ref{thm:full-spectrum}), the dimension theorems
(Theorems~\ref{thm:finite}, \ref{thm:upper} and~\ref{conj:weight}), the exact
gauges (Theorem~\ref{thm:gauge} and Corollaries~\ref{cor:exact-gauge}
and~\ref{cor:liminf-gauges}) and the trichotomy
(Theorem~\ref{thm:trichotomy}).

For the weight $\alpha=0$, the dimension statement of Theorem~\ref{conj:weight}
also follows from Corollary~8 of Bugeaud, Liao and Rams \cite{BLR}, which applies
because $Q_{j+1}/Q_j\to\infty$ (in particular $\dimH\{x:\|Q_jx\|\to0\}=1$). For
$0<\alpha<1$ the logarithmic growth condition~(1.6) of their Theorem~7 fails,
since $\log Q_{j+1}/\log Q_j\to1$. As the maps $x\mapsto Q_jx$ are linear,
the lower bound can instead be read from the homogeneous Moran bounds
of~\cite{FWW}; we give a direct covering upper bound and a shifted-grid Moran
construction. Mass transference methods for limsup sets \cite{BV} do not apply
to the eventual (liminf-type) conditions considered here.

Regarding gauges, Cabrelli, Mendivil, Molter and Shonkwiler
\cite[\S5, Theorem~3]{CMMS} construct for a Cantor set a gauge giving it
positive \emph{finite} measure, whereas Moreira, Erazo and Angelini \cite{MEA}
show that the classical Lagrange level at value $3$ has zero or
non-$\sigma$-finite measure for every dimension function. Here the recurrence
supplies the gauge $h_k$, and the level $\{\Lam=k/4\}$ has \emph{infinite but
$\sigma$-finite} $\mathcal H^{h_k}$-measure.

The phenomenon of a spectral gap with zero-dimensional sub-threshold levels is
classical for the Lagrange spectrum, whose part below $3$ is countable while
Hall's ray gives a half-line; Moreira \cite{Moreira} studies the dimension of
the Lagrange and Markov spectra below a threshold, and Lima and
Moreira~\cite{LimaMoreira} prove phase transitions for dynamical Markov and
Lagrange spectra of horseshoes. Our level sets live in the space of inputs
$x\in\R$, not in the spectrum. For Cantor series expansions, Cheng and
Zhang~\cite{CZ} determine Hausdorff dimensions of related limsup sets; spectra
and exact gauges are not considered there.

\subsection{Formal verification}

The results of this paper have also been formalized in the Lean~4 proof
assistant with the Mathlib library, using only Lean's standard axioms. The
formal development and its axiom audit are described in the supplementary
material. The proofs below are complete and do not depend on the
formalization.

\subsection{Organization}

Section~\ref{sec:class} introduces the digit coding and proves the zero-level
classification and the initial gap. Section~\ref{sec:spectrum} proves the
quantization of the spectrum. Section~\ref{sec:finite-levels} shows that all
finite levels have dimension $0$, and Section~\ref{sec:weights} treats the
subcritical weights. Section~\ref{sec:gauge} contains the exact gauges and the
trichotomy.
